% =================================================
% Ученический рабочий лист. Урок 1
% =================================================

\worksheettitle
  {Цифры, числа и разряды}
  {Урок 1 из 3. Обозначение натуральных чисел}

\studentnameblock

\begin{checkpointbox}[Входная диагностика: 4 минуты]
  Выполните задания без подробных объяснений.
  \begin{enumerate}[label=\textbf{\arabic*.},itemsep=0.38em]
    \item Сколько цифр в записи числа \(405\)?
      Ответ: \answerblank[15mm].
    \item Какая цифра стоит в разряде десятков числа \(684\)?
      Ответ: \answerblank[15mm].
    \item \(307=\answerblank[43mm].\)
    \item \(4\,000+50+2=\answerblank[25mm].\)
  \end{enumerate}
\end{checkpointbox}

\section{Из чего записывают числа}

Заполните пропуски:
\[
  0,\ 1,\ 2,\ \answerblank[8mm],\ 4,\ 5,\ \answerblank[8mm],\ 7,\ 8,\ \answerblank[8mm].
\]

\begin{fillbox}[Исследуем записи]
  \begin{center}
    \renewcommand{\arraystretch}{1.5}
    \begin{tabularx}{0.92\textwidth}{@{}>{\centering\arraybackslash}p{0.2\textwidth}>{\centering\arraybackslash}p{0.28\textwidth}X@{}}
      \toprule
      \rowcolor{TableHeader}
      Запись числа & Сколько цифр? & Цифры по порядку \\
      \midrule
      \(8\) & \answerblank[12mm] & \answerblank[30mm] \\
      \(47\) & \answerblank[12mm] & \answerblank[30mm] \\
      \(505\) & \answerblank[12mm] & \answerblank[40mm] \\
      \bottomrule
    \end{tabularx}
  \end{center}
\end{fillbox}

Дополните определения.

\begin{itemize}
  \item \textbf{Цифра} --- это \answerblank[70mm].
  \item \textbf{Число} --- это \answerblank[70mm].
\end{itemize}

\begin{thinkbox}[Однозначное или многозначное?]
  Подчеркните однозначные числа одной чертой, а многозначные --- двумя:
  \[
    7,\qquad 12,\qquad 304,\qquad 9,\qquad 6\,000.
  \]
\end{thinkbox}

\section{Одна цифра --- разные значения}

Заполните таблицу.

\begin{center}
  \renewcommand{\arraystretch}{1.5}
  \begin{tabularx}{0.94\textwidth}{@{}>{\centering\arraybackslash}p{0.16\textwidth}>{\centering\arraybackslash}p{0.32\textwidth}X@{}}
    \toprule
    \rowcolor{TableHeader}
    Число & Разряд цифры \(6\) & Значение цифры \(6\) \\
    \midrule
    \(6\) & единицы & \(6\) \\
    \(60\) & \answerblank[28mm] & \answerblank[22mm] \\
    \(600\) & \answerblank[28mm] & \answerblank[22mm] \\
    \(6\,000\) & \answerblank[28mm] & \answerblank[22mm] \\
    \bottomrule
  \end{tabularx}
\end{center}

\begin{reflectionbox}[Сформулируйте вывод]
  Значение цифры зависит от
  \answerblank[85mm].
\end{reflectionbox}

\section{Разряды многозначного числа}

Рассмотрите число \(72\,649\).

\placevalueworksheetfigure

\begin{fillbox}[Прочитайте разрядную запись]
  \begin{enumerate}
    \item Цифра десятков тысяч: \answerblank[14mm].
    \item Цифра единиц тысяч: \answerblank[14mm].
    \item Цифра сотен: \answerblank[14mm].
    \item Цифра десятков: \answerblank[14mm].
    \item Цифра единиц: \answerblank[14mm].
  \end{enumerate}
\end{fillbox}

Каково значение каждой цифры?
\[
  7\longrightarrow\answerblank[22mm],\qquad
  2\longrightarrow\answerblank[22mm],\qquad
  6\longrightarrow\answerblank[18mm],
\]
\[
  4\longrightarrow\answerblank[18mm],\qquad
  9\longrightarrow\answerblank[18mm].
\]

\section{Ноль сохраняет место разряда}

Рассмотрите число \(40\,508\).

\zeroplacefigure

\begin{checkpointbox}[Что показывают нули?]
  \begin{enumerate}
    \item Цифра в разряде тысяч равна \answerblank[12mm].
    \item Цифра в разряде десятков равна \answerblank[12mm].
    \item Если удалить оба нуля, получится число \answerblank[20mm].
    \item Равно ли оно числу \(40\,508\)?
      \choicebox\ да \qquad \choicebox\ нет.
  \end{enumerate}
\end{checkpointbox}

Дополните вывод:

\begin{reflectionbox}
  Ноль внутри записи числа показывает, что в соответствующем разряде
  \answerblank[75mm].
\end{reflectionbox}

\section{Разложение числа по разрядам}

\expandedformfigure

Заполните пропуски:
\[
  47\,305=
  \answerblank[20mm]+\answerblank[20mm]+\answerblank[18mm]+\answerblank[15mm].
\]

\begin{fillbox}[Разложите числа]
  \begin{enumerate}
    \item \(63\,204=\answerblank[90mm].\)
    \item \(8\,090=\answerblank[90mm].\)
  \end{enumerate}
\end{fillbox}

\begin{fillbox}[Соберите числа]
  \begin{enumerate}
    \item \(70\,000+800+6=\answerblank[25mm].\)
    \item \(5\,000+40+3=\answerblank[25mm].\)
    \item \(40\,000+2\,000+50+9=\answerblank[25mm].\)
  \end{enumerate}
\end{fillbox}

\begin{warningbox}[Найдите и исправьте ошибку]
  Ученик записал:
  \[
    50\,204=50\,000+2\,000+4.
  \]
  Ошибка состоит в том, что
  \writinglines{2}
  Верная запись:
  \[
    50\,204=\answerblank[90mm].
  \]
\end{warningbox}

\section{Самостоятельная проверка}

\begin{practicebox}[Выполните без подсказок]
  \begin{enumerate}[label=\textbf{\arabic*.},itemsep=0.38em]
    \item Сколько цифр в записи числа \(70\,604\)?
      Ответ: \answerblank[18mm].
    \item Какая цифра стоит в разряде сотен числа \(47\,382\)?
      Ответ: \answerblank[18mm].
    \item Каково значение цифры \(6\) в числе \(26\,415\)?
      Ответ: \answerblank[23mm].
    \item Разложите число \(90\,307\) по разрядам:
      \answerblank[64mm].
    \item Соберите число: \(30\,000+4\,000+200+7=\answerblank[22mm].\)
    \item Запишите число, в котором \(4\) десятка тысяч, \(7\) сотен и
      \(2\) единицы: \answerblank[22mm].
  \end{enumerate}
\end{practicebox}

\begin{checkpointbox}[Выходной билет]
  \begin{enumerate}
    \item Число \(608\) записано \answerblank[12mm] цифрами.
    \item \(30\,000+400+8=\answerblank[22mm].\)
    \item Цифра десятков числа \(61\,507\) равна \answerblank[12mm].
  \end{enumerate}
\end{checkpointbox}

\begin{reflectionbox}[Оцените себя]
  \choicebox\ могу объяснить другому;
  \qquad
  \choicebox\ умею по образцу;
  \qquad
  \choicebox\ нужна тренировка.
\end{reflectionbox}
